% GENERATED FILE. Edit canonical lessons, then run npm run generate:books.
% canonical-source-hash: fnv1a64:ac63d2991562f0da
% canonical-lessons: TE-C385-angikarincu, TE-C385-angikaram, TE-C385-ekibhavincu, TE-C385-sammatincu, TE-C385-maddatu

\chapter{Agreeing}
\label{ch:te-agreeing}
\hlchaptermodality{385}

\begin{chapteropening}
\textbf{By the end of this chapter:} \emph{I can agree, accept and give consent.}

Complete the last lesson of chapter 385: I can agree, accept and give consent.
\end{chapteropening}

\section[\texorpdfstring{\te{అంగీకరించు} (aṅgīkariñcu)}{అంగీకరించు (aṅgīkariñcu)}]{\te{అంగీకరించు} (aṅgīkariñcu) — to accept; to agree to}

\label{lesson:TE-C385-angikarincu}



\begin{quote}
Before the new one: say the Telugu for to recommend, then the Telugu for a point.
\end{quote}

\subsection*{You'll want to know: \te{అంగీకరించు}}
\textbf{\te{అంగీకరించు}} — \emph{aṅgīkariñcu} — \textquotedblleft{}to accept; to agree to\textquotedblright{}.

\begin{quote}
\te{కంపెనీ} \te{మా} \te{షరతులు} \te{అంగీకరించింది}.
\end{quote}

\begin{quote}
\emph{kampenī mā ṣaratulu aṅgīkariñcindi.}
\end{quote}

\begin{quote}
The company accepted our conditions.
\end{quote}

Past (I …): \te{అంగీకరించాను}, \emph{aṅgīkariñcānu}.

Formal beside the everyday \te{ఒప్పుకో}.

\begin{cousinweb}
\textbf{\te{అంగీకరించు}} (\emph{aṅgīkariñcu}) is built on Sanskrit \emph{aṅgīkāra}, \textquotedblleft{}agreement, a promise\textquotedblright{} (Monier-Williams), with the Telugu verb ending -\te{ఇంచు}.
\end{cousinweb}

\begin{grammarlens}[title={What you've built}]
One more word for giving reasons and opinions.
\end{grammarlens}

\subsection*{Your turn}
\begin{itemize}
\raggedright
  \item \emph{Say it:} \emph{aṅgīkariñcu}
  \item \emph{Say it:} \emph{aṅgīkariñcu}, once more
  \item \emph{Say it:} \emph{kampenī mā ṣaratulu aṅgīkariñcindi.} — The company accepted our conditions.
  \item \emph{Recall:} say the Telugu for to recommend, then the Telugu for a point, then say \emph{aṅgīkariñcu} again
\end{itemize}

\begin{tcolorbox}[breakable,colback=teal!4,colframe=teal!35!black,title={Before you move on}]
What does \te{అంగీకరించు} mean? (\textquotedblleft{}To accept; to agree to\textquotedblright{}.) Say it once more.
\end{tcolorbox}
\section[\texorpdfstring{\te{అంగీకారం} (aṅgīkāraṁ)}{అంగీకారం (aṅgīkāraṁ)}]{\te{అంగీకారం} (aṅgīkāraṁ) — consent; acceptance}

\label{lesson:TE-C385-angikaram}



\begin{quote}
Before the new one: say the Telugu for a point, then the Telugu for to accept.
\end{quote}

\subsection*{You'll want to know: \te{అంగీకారం}}
\textbf{\te{అంగీకారం}} — \emph{aṅgīkāraṁ} — \textquotedblleft{}consent; acceptance\textquotedblright{}.

\begin{quote}
\te{తల్లిదండ్రుల} \te{అంగీకారంతో} \te{వాళ్ళ} \te{పెళ్ళి} \te{జరిగింది}.
\end{quote}

\begin{quote}
\emph{tallidaṇḍrula aṅgīkārantō vāḷḷa peḷḷi jarigindi.}
\end{quote}

\begin{quote}
Their wedding took place with their parents’ consent.
\end{quote}

The noun of \te{అంగీకరించు}: \te{అంగీకారం} \te{తెలుపు}, to give one’s consent.

\begin{cousinweb}
\textbf{\te{అంగీకారం}} (\emph{aṅgīkāraṁ}) is borrowed from Sanskrit \emph{aṅgīkāra}, \textquotedblleft{}agreement, a promise\textquotedblright{} (Monier-Williams); its verb is \te{అంగీకరించు}.
\end{cousinweb}

\begin{grammarlens}[title={What you've built}]
One more word for giving reasons and opinions.
\end{grammarlens}

\subsection*{Your turn}
\begin{itemize}
\raggedright
  \item \emph{Say it:} \emph{aṅgīkāraṁ}
  \item \emph{Say it:} \emph{aṅgīkāraṁ}, once more
  \item \emph{Say it:} \emph{tallidaṇḍrula aṅgīkārantō vāḷḷa peḷḷi jarigindi.} — Their wedding took place with their parents’ consent.
  \item \emph{Recall:} say the Telugu for a point, then the Telugu for to accept, then say \emph{aṅgīkāraṁ} again
\end{itemize}

\begin{tcolorbox}[breakable,colback=teal!4,colframe=teal!35!black,title={Before you move on}]
What does \te{అంగీకారం} mean? (\textquotedblleft{}Consent; acceptance\textquotedblright{}.) Say it once more.
\end{tcolorbox}
\section[\texorpdfstring{\te{ఏకీభవించు} (ēkībhaviñcu)}{ఏకీభవించు (ēkībhaviñcu)}]{\te{ఏకీభవించు} (ēkībhaviñcu) — to agree with}

\label{lesson:TE-C385-ekibhavincu}



\begin{quote}
Before the new one: say the Telugu for to accept, then the Telugu for consent.
\end{quote}

\subsection*{You'll want to know: \te{ఏకీభవించు}}
\textbf{\te{ఏకీభవించు}} — \emph{ēkībhaviñcu} — \textquotedblleft{}to agree with\textquotedblright{}.

\begin{quote}
\te{నేను} \te{మీ} \te{అభిప్రాయంతో} \te{ఏకీభవిస్తాను}.
\end{quote}

\begin{quote}
\emph{nēnu mī abhiprāyantō ēkībhavistānu.}
\end{quote}

\begin{quote}
I agree with your opinion.
\end{quote}

Past (I …): \te{ఏకీభవించాను}, \emph{ēkībhaviñcānu}.

The person or view you agree with takes -\te{తో}: \te{మీతో} \te{ఏకీభవిస్తాను}.

\begin{cousinweb}
\textbf{\te{ఏకీభవించు}} (\emph{ēkībhaviñcu}) is built on Sanskrit \emph{ekībhāva}, \textquotedblleft{}the becoming one, coalition\textquotedblright{} (Monier-Williams), from \emph{eka}, \textquotedblleft{}one\textquotedblright{}, and \emph{bhāva}, \textquotedblleft{}becoming\textquotedblright{}, with the Telugu verb ending -\te{ఇంచు}.
\end{cousinweb}

\begin{grammarlens}[title={What you've built}]
One more word for giving reasons and opinions.
\end{grammarlens}

\subsection*{Your turn}
\begin{itemize}
\raggedright
  \item \emph{Say it:} \emph{ēkībhaviñcu}
  \item \emph{Say it:} \emph{ēkībhaviñcu}, once more
  \item \emph{Say it:} \emph{nēnu mī abhiprāyantō ēkībhavistānu.} — I agree with your opinion.
  \item \emph{Recall:} say the Telugu for to accept, then the Telugu for consent, then say \emph{ēkībhaviñcu} again
\end{itemize}

\begin{tcolorbox}[breakable,colback=teal!4,colframe=teal!35!black,title={Before you move on}]
What does \te{ఏకీభవించు} mean? (\textquotedblleft{}To agree with\textquotedblright{}.) Say it once more.
\end{tcolorbox}
\section[\texorpdfstring{\te{సమ్మతించు} (sammatiñcu)}{సమ్మతించు (sammatiñcu)}]{\te{సమ్మతించు} (sammatiñcu) — to consent}

\label{lesson:TE-C385-sammatincu}



\begin{quote}
Before the new one: say the Telugu for consent, then the Telugu for to agree with.
\end{quote}

\subsection*{You'll want to know: \te{సమ్మతించు}}
\textbf{\te{సమ్మతించు}} — \emph{sammatiñcu} — \textquotedblleft{}to consent\textquotedblright{}.

\begin{quote}
\te{మా} \te{నాన్న} \te{ఆ} \te{ప్రయాణానికి} \te{సమ్మతించాడు}.
\end{quote}

\begin{quote}
\emph{mā nānna ā prayāṇāniki sammatiñcāḍu.}
\end{quote}

\begin{quote}
My father consented to that trip.
\end{quote}

Past (I …): \te{సమ్మతించాను}, \emph{sammatiñcānu}.

What you consent to takes -\te{కి}: \te{దానికి} \te{సమ్మతించు}. The noun is \te{సమ్మతి}, consent.

\begin{cousinweb}
\textbf{\te{సమ్మతించు}} (\emph{sammatiñcu}) is built on Sanskrit \emph{sammati}, \textquotedblleft{}sameness of opinion, agreement, approval\textquotedblright{} (Monier-Williams), with the Telugu verb ending -\te{ఇంచు}.
\end{cousinweb}

\begin{grammarlens}[title={What you've built}]
One more word for giving reasons and opinions.
\end{grammarlens}

\subsection*{Your turn}
\begin{itemize}
\raggedright
  \item \emph{Say it:} \emph{sammatiñcu}
  \item \emph{Say it:} \emph{sammatiñcu}, once more
  \item \emph{Say it:} \emph{mā nānna ā prayāṇāniki sammatiñcāḍu.} — My father consented to that trip.
  \item \emph{Recall:} say the Telugu for consent, then the Telugu for to agree with, then say \emph{sammatiñcu} again
\end{itemize}

\begin{tcolorbox}[breakable,colback=teal!4,colframe=teal!35!black,title={Before you move on}]
What does \te{సమ్మతించు} mean? (\textquotedblleft{}To consent\textquotedblright{}.) Say it once more.
\end{tcolorbox}
\section[\texorpdfstring{\te{మద్దతు} (maddatu)}{మద్దతు (maddatu)}]{\te{మద్దతు} (maddatu) — backing; support}

\label{lesson:TE-C385-maddatu}



\begin{quote}
Before the new one: say the Telugu for to agree with, then the Telugu for to consent.
\end{quote}

\subsection*{You'll want to know: \te{మద్దతు}}
\textbf{\te{మద్దతు}} — \emph{maddatu} — \textquotedblleft{}backing; support\textquotedblright{}.

\begin{quote}
\te{ఆ} \te{ఆలోచనకు} \te{చాలా} \te{మంది} \te{మద్దతు} \te{ఇచ్చారు}.
\end{quote}

\begin{quote}
\emph{ā ālōcanaku cālā mandi maddatu iccāru.}
\end{quote}

\begin{quote}
Many people backed that idea.
\end{quote}

\te{మద్దతు} \te{ఇచ్చు}, to give support; the person or idea backed takes -\te{కి}.

\begin{cousinweb}
\textbf{\te{మద్దతు}} (\emph{maddatu}) is usually traced to Hindustani \emph{madad}, \textquotedblleft{}help\textquotedblright{}, a word of Arabic origin; the Telugu form and its sense of backing a cause are its own.
\end{cousinweb}

\begin{grammarlens}[title={What you've built}]
One more word for giving reasons and opinions.
\end{grammarlens}

\subsection*{Your turn}
\begin{itemize}
\raggedright
  \item \emph{Say it:} \emph{maddatu}
  \item \emph{Say it:} \emph{maddatu}, once more
  \item \emph{Say it:} \emph{ā ālōcanaku cālā mandi maddatu iccāru.} — Many people backed that idea.
  \item \emph{Recall:} say the Telugu for to agree with, then the Telugu for to consent, then say \emph{maddatu} again
\end{itemize}

\begin{tcolorbox}[breakable,colback=teal!4,colframe=teal!35!black,title={Before you move on}]
What does \te{మద్దతు} mean? (\textquotedblleft{}Backing; support\textquotedblright{}.) Say it once more.
\end{tcolorbox}
